% p406_1; source PDF page 406; manually verified editable table.
\begin{tikzpicture}[x=1cm,y=-1cm,inner sep=0pt,outer sep=0pt]
\path[use as bounding box] (0,0) rectangle (16.000,9.225);
\draw[line width=.4pt] (0.000,0.000) rectangle (1.900,0.915);
\node[align=center,text width=1.640cm,font={\linespread{1}\fontsize{10}{13.50}\selectfont\bfseries}] at (0.950,0.458) {罪名};
\draw[line width=.4pt] (1.900,0.000) rectangle (5.100,0.915);
\node[align=center,text width=2.940cm,font={\linespread{1}\fontsize{10}{13.50}\selectfont\bfseries}] at (3.500,0.458) {挪用资金罪};
\draw[line width=.4pt] (5.100,0.000) rectangle (7.900,0.915);
\node[align=center,text width=2.540cm,font={\linespread{1}\fontsize{10}{13.50}\selectfont\bfseries}] at (6.500,0.458) {挪用公款罪};
\draw[line width=.4pt] (7.900,0.000) rectangle (16.000,0.915);
\node[align=center,text width=7.840cm,font={\linespread{1}\fontsize{10}{13.50}\selectfont\bfseries}] at (11.950,0.458) {挪用特定款物罪};
\draw[line width=.4pt] (0.000,0.915) rectangle (1.900,1.980);
\node[align=center,text width=1.640cm,font={\linespread{1}\fontsize{10}{13.50}\selectfont}] at (0.950,1.448) {行为主体};
\draw[line width=.4pt] (1.900,0.915) rectangle (5.100,1.980);
\node[align=center,text width=2.940cm,font={\linespread{1}\fontsize{10}{13.50}\selectfont}] at (3.500,1.448) {非国家工作人员};
\draw[line width=.4pt] (5.100,0.915) rectangle (7.900,1.980);
\node[align=center,text width=2.540cm,font={\linespread{1}\fontsize{10}{13.50}\selectfont}] at (6.500,1.448) {国家工作人员};
\draw[line width=.4pt] (7.900,0.915) rectangle (16.000,1.980);
\node[align=left,text width=7.840cm,font={\linespread{1}\fontsize{10}{13.50}\selectfont}] at (11.950,1.448) {一般主体，包括国家工作人员和非国家工作人员①};
\draw[line width=.4pt] (0.000,1.980) rectangle (1.900,3.595);
\node[align=center,text width=1.640cm,font={\linespread{1}\fontsize{10}{13.50}\selectfont}] at (0.950,2.787) {行为对象};
\draw[line width=.4pt] (1.900,1.980) rectangle (5.100,3.595);
\node[align=center,text width=2.940cm,font={\linespread{1}\fontsize{10}{13.50}\selectfont}] at (3.500,2.787) {非国有单位的资金};
\draw[line width=.4pt] (5.100,1.980) rectangle (7.900,3.595);
\node[align=center,text width=2.540cm,font={\linespread{1}\fontsize{10}{13.50}\selectfont}] at (6.500,2.787) {公家的资金};
\draw[line width=.4pt] (7.900,1.980) rectangle (16.000,3.595);
\node[align=left,text width=7.840cm,font={\linespread{1}\fontsize{10}{13.50}\selectfont}] at (11.950,2.787) {特定款物，指专用于抢险、救灾、救济等用途的款（资金）和物（物品）};
\draw[line width=.4pt] (0.000,3.595) rectangle (1.900,5.210);
\node[align=center,text width=1.640cm,font={\linespread{1}\fontsize{10}{13.50}\selectfont}] at (0.950,4.402) {行为方式};
\draw[line width=.4pt] (1.900,3.595) rectangle (5.100,5.210);
\node[align=left,text width=2.940cm,font={\linespread{1}\fontsize{10}{13.50}\selectfont}] at (3.500,4.402) {挪作私用，有归还意图};
\draw[line width=.4pt] (5.100,3.595) rectangle (7.900,5.210);
\node[align=left,text width=2.540cm,font={\linespread{1}\fontsize{10}{13.50}\selectfont}] at (6.500,4.402) {挪作私用，有归还意图};
\draw[line width=.4pt] (7.900,3.595) rectangle (16.000,5.210);
\node[align=left,text width=7.840cm,font={\linespread{1}\fontsize{10}{13.50}\selectfont}] at (11.950,4.402) {将A公用用途变成B公用用途，有无归还意图，在所不问};
\draw[line width=.4pt] (0.000,5.210) rectangle (1.900,9.225);
\node[align=center,text width=1.640cm,font={\linespread{1}\fontsize{10}{13.50}\selectfont}] at (0.950,7.217) {要点};
\draw[line width=.4pt] (1.900,5.210) rectangle (16.000,9.225);
\node[align=left,text width=13.840cm,font={\linespread{1}\fontsize{10}{13.50}\selectfont}] at (8.950,7.217) {1.将特定款物，挪作私用，定挪用公款罪。其中的“特定物”定挪用公款罪，属于特殊规定，因为“物”本不属于“款”。\\{}2.失业保险基金、下岗职工基本生活保障资金，属于救济类特定款物。根据行为方式不同，可定挪用特定款物罪或挪用公款罪。\\{}3.退休职工养老保险金，不属于特定款物。根据行为主体身份不同，可定挪用资金罪或挪用公款罪};
\end{tikzpicture}
